\documentclass[a4paper]{arrowhead}

\usepackage[yyyymmdd]{datetime}
\usepackage{etoolbox}
\usepackage[T1]{fontenc}
\usepackage[utf8]{inputenc}
\usepackage{multirow}
\usepackage{longtable}
\usepackage{amssymb}
\usepackage{float}
\usepackage{tikz}
\usetikzlibrary{positioning,arrows.meta,fit,calc,shapes.misc,shapes.geometric}

\renewcommand{\dateseparator}{-}
\setlength{\emergencystretch}{3em}
\lstset{basicstyle=\ttfamily\footnotesize,breaklines=true,columns=fullflexible,frame=single,rulecolor=\color{gray!50},backgroundcolor=\color{gray!6}}
\tikzset{
  ahsys/.style={draw=ArrowheadBlue,very thick,fill=ArrowheadBlue!6,rounded corners=2pt,minimum width=3.2cm,minimum height=1.1cm,align=center,font=\small\bfseries},
  ahext/.style={draw=ArrowheadGrey,thick,dashed,fill=gray!5,rounded corners=2pt,minimum width=3.2cm,minimum height=1.0cm,align=center,font=\small},
  ahsrv/.style={draw=ArrowheadPurple,thick,fill=ArrowheadPurple!8,rounded corners=6pt,minimum height=0.55cm,inner xsep=4pt,align=center,font=\scriptsize\ttfamily},
  ahcsm/.style={draw=ArrowheadTeal,thick,fill=ArrowheadTeal!8,rounded corners=6pt,minimum height=0.55cm,inner xsep=4pt,align=center,font=\scriptsize\ttfamily},
  ahdb/.style={draw=ArrowheadOrange,thick,fill=ArrowheadOrange!8,cylinder,shape border rotate=90,aspect=0.25,minimum height=0.9cm,align=center,font=\scriptsize},
  ahflow/.style={-{Stealth[length=2.2mm]},thick,ArrowheadGrey},
  ahlabel/.style={font=\scriptsize,fill=white,inner sep=1pt,align=center},
}

\newcommand{\fparam}[1]{\textit{\textcolor{ArrowheadBlue}{#1}}}
%% Special references
\newcommand{\mref}[1]{{\textcolor{ArrowheadPurple}{\hyperref[sec:model:#1]{#1}}}}
\newcommand{\oref}[1]{{\textcolor{ArrowheadBlue}{\hyperref[sec:operations:#1]{#1}}}}
\newcommand{\pdef}[1]{{\textcolor{ArrowheadGrey}{#1\label{sec:model:primitives:#1}\label{sec:model:primitives:#1s}}}}
\newcommand{\pref}[1]{{\textcolor{ArrowheadGrey}{\hyperref[sec:model:primitives:#1]{#1}}}}

\newrobustcmd\osubsection[5]{
  \addtocounter{subsection}{1}
  \addcontentsline{toc}{subsection}{\protect\numberline{\thesubsection}operation \textcolor{ArrowheadBlue}{#1}}
  \subsection*{
    \thesubsection\quad
    #2 \textcolor{ArrowheadPurple}{#3} \\
    \small
    \hspace*{0.075\textwidth}\begin{minipage}{0.1\textwidth}
      \vspace*{1mm}
      Operation: \\
      \notblank{#4}{Input: \\}{}
      \notblank{#5}{Output: \\}{}
    \end{minipage}
    \begin{minipage}{0.825\textwidth}
      \vspace*{1mm}
      \textcolor{ArrowheadBlue}{#1} \\
      \notblank{#4}{\mref{#4} \\}{}
      \notblank{#5}{\mref{#5} \\}{}
    \end{minipage}
  }
  \label{sec:operations:#1}
}
\newrobustcmd\msubsection[2]{
  \addtocounter{subsection}{1}
  \addcontentsline{toc}{subsection}{\protect\numberline{\thesubsection}#1 \textcolor{ArrowheadPurple}{#2}}
  \subsection*{\thesubsection\quad#1 \textcolor{ArrowheadPurple}{#2}}
  \label{sec:model:#2} \label{sec:model:#2s}
}

\begin{document}

%% Arrowhead Document Properties
\ArrowheadTitle{workflowRecipe HTTP/TLS/JSON}
\ArrowheadServiceID{workflowRecipe}
\ArrowheadType{Interface Design Description}
\ArrowheadTypeShort{IDD}
\ArrowheadVersion{5.4.0}
\ArrowheadDate{2026-11-23}
\ArrowheadAuthor{Documentation maintenance}
\ArrowheadStatus{REVIEW DRAFT}
\ArrowheadContact{Repository maintainers}
\ArrowheadFooter{\href{www.arrowhead.eu}{www.arrowhead.eu}}
\ArrowheadSetup
%%

%% Front Page
\begin{center}
  \vspace*{1cm}
  \huge{\arrowtitle}

  \vspace*{0.2cm}
  \LARGE{\arrowtype}
  \vspace*{1cm}

  \Large{Service ID: \textit{``\arrowid''}}
  \vspace*{\fill}

  

  \vspace*{1cm}
  \vspace*{\fill}

  % Front Page Abstract
  \begin{abstract}
    This document describes the HTTP/TLS/JSON interface of the \texttt{workflowRecipe} service of the Eclipse Arrowhead workflow support systems, service version 1.0.0, documentation edition 5.4.0. It realises the abstract operations of the \texttt{workflowRecipe} SD.
  \end{abstract}

  \vspace*{1cm}
\end{center}

\newpage
%%

%% Table of Contents
\tableofcontents
\newpage
%%

\section{Overview}
\label{sec:overview}
This document describes how the \texttt{workflowRecipe} service, which validates, registers, lists, retires and compiles workflow recipes, is realised with HTTP~\cite{rfc9110}
over TLS~\cite{rfc8446} and JSON~\cite{rfc8259} payloads. Request and response bodies use content type
\texttt{application/json}; JSON field names are case-sensitive and fields not listed in this document must not be sent
(the recipe decoder rejects unknown fields). Optional fields may be omitted; absent and \texttt{null} are equivalent.

The service is offered by the WorkflowManager on its Arrowhead connector (default port 8482) with base path \texttt{/\allowbreak{}workflowmanager/\allowbreak{}recipe}.
The ServiceRegistry entry lists every operation with its HTTP method and path relative to the base path.

The rest of this document is organized as follows. Section~\ref{sec:operations} specifies each operation,
Section~\ref{sec:model} the JSON data models, Section~\ref{sec:security} the security requirements,
Section~\ref{sec:errors} the error representation and Section~\ref{sec:implementations} the implementations.

\newpage
\section{Service Operations}
\label{sec:operations}

{\small
\begin{longtable}{|p{3.0cm}|p{1.5cm}|p{5.3cm}|p{3.0cm}|p{3.0cm}|}

\hline
\rowcolor{gray!33} Operation & Method & Path & Input & Output \\ \hline
\endhead
\oref{validate} & POST & \texttt{/\allowbreak{}workflowmanager/\allowbreak{}recipe/\allowbreak{}validate} & \mref{Recipe} & \mref{ValidationReport} \\ \hline
\oref{register} & POST & \texttt{/\allowbreak{}workflowmanager/\allowbreak{}recipe/\allowbreak{}register} & \mref{Recipe} & \mref{RecipeRegistration} \\ \hline
\oref{query} & GET & \texttt{/\allowbreak{}workflowmanager/\allowbreak{}recipe/\allowbreak{}recipes} & \mref{RecipeFilter} & \mref{RecipeList} \\ \hline
\oref{get} & GET & \texttt{/\allowbreak{}workflowmanager/\allowbreak{}recipe/\allowbreak{}recipes/\allowbreak{}\{name\}/\allowbreak{}\{version\}} & \mref{RecipeKey} & \mref{RecipeRecord} \\ \hline
\oref{retire} & DELETE & \texttt{/\allowbreak{}workflowmanager/\allowbreak{}recipe/\allowbreak{}recipes/\allowbreak{}\{name\}/\allowbreak{}\{version\}} & \mref{RecipeKey} & \mref{RecipeRecord} \\ \hline
\oref{compile} & POST & \texttt{/\allowbreak{}workflowmanager/\allowbreak{}recipe/\allowbreak{}compile} & \mref{RecipeKey} & \mref{ChoreographerPlan} \\ \hline
\end{longtable}
}

\osubsection{validate}{POST}{/workflowmanager/recipe/validate}{Recipe}{ValidationReport}
Validates a recipe without storing it. Every finding is reported; the verdict is the worst outcome.

{\small
\begin{longtable}{|p{2.0cm}|p{13.8cm}|}

\hline
\rowcolor{gray!33} Status & When \\ \hline
\endhead
Success & 200 OK with the report, whatever the verdict. \\ \hline
400 & The body is not a recipe (unknown field, wrong type). \\ \hline
401 & No valid identity token, or the token is bound to another certificate (\texttt{AUTH}). \\ \hline
403 & The consumer has no ConsumerAuthorization grant for this operation (\texttt{FORBIDDEN}). \\ \hline
\end{longtable}
}


\begin{lstlisting}[caption={Example request body}]
{
  "name": "temperature-branching",
  "version": "1.0.0",
  "firstAction": "control",
  "parameters": [ { "name": "threshold", "type": "number" } ],
  "actions": [ {
    "name": "control",
    "steps": [
      { "name": "read", "service": { "serviceDefinition": "temperature-read" },
        "next": ["evaluate"] },
      { "name": "evaluate", "service": { "serviceDefinition": "temperature-evaluate" },
        "next": ["heatOn", "heatOff"], "dep": ["read"],
        "input": { "temperature": "${result.read.temperature}",
                   "threshold": "${param.threshold}" },
        "guards": { "heatOn": "${result.evaluate.heat} == true",
                    "heatOff": "!${result.evaluate.heat}" } },
      { "name": "heatOn", "service": { "serviceDefinition": "heating-switch" },
        "next": ["confirm"], "input": { "on": true } },
      { "name": "heatOff", "service": { "serviceDefinition": "heating-switch" },
        "next": ["confirm"], "input": { "on": false } },
      { "name": "confirm", "service": { "serviceDefinition": "temperature-read" },
        "dep": ["heatOn", "heatOff"] }
    ] } ]
}
\end{lstlisting}


\begin{lstlisting}[caption={Example response body}]
{ "recipe": "temperature-branching@1.0.0", "verdict": "ACCEPT",
  "findings": [ { "code": "G04", "outcome": "ACCEPT", "basis": "x-guard",
    "location": "recipe", "message": "Recipe uses guarded transitions (extension x-guard)" } ] }
\end{lstlisting}

\osubsection{register}{POST}{/workflowmanager/recipe/register}{Recipe}{RecipeRegistration}
Validates the recipe and stores it when the verdict is ACCEPT. The registering system is recorded as \texttt{registeredBy}.

{\small
\begin{longtable}{|p{2.0cm}|p{13.8cm}|}

\hline
\rowcolor{gray!33} Status & When \\ \hline
\endhead
Success & 201 Created for a new version; 200 OK when identical content is already registered. \\ \hline
400 & The verdict is not ACCEPT; the body is a \mref{RecipeRegistration} with \texttt{recipe = null} and the report. \\ \hline
409 & The name and version exist with different content. \\ \hline
401 & No valid identity token, or the token is bound to another certificate (\texttt{AUTH}). \\ \hline
403 & The consumer has no ConsumerAuthorization grant for this operation (\texttt{FORBIDDEN}). \\ \hline
\end{longtable}
}


\begin{lstlisting}[caption={Example request body}]
{
  "name": "temperature-branching",
  "version": "1.0.0",
  "firstAction": "control",
  "parameters": [ { "name": "threshold", "type": "number" } ],
  "actions": [ {
    "name": "control",
    "steps": [
      { "name": "read", "service": { "serviceDefinition": "temperature-read" },
        "next": ["evaluate"] },
      { "name": "evaluate", "service": { "serviceDefinition": "temperature-evaluate" },
        "next": ["heatOn", "heatOff"], "dep": ["read"],
        "input": { "temperature": "${result.read.temperature}",
                   "threshold": "${param.threshold}" },
        "guards": { "heatOn": "${result.evaluate.heat} == true",
                    "heatOff": "!${result.evaluate.heat}" } },
      { "name": "heatOn", "service": { "serviceDefinition": "heating-switch" },
        "next": ["confirm"], "input": { "on": true } },
      { "name": "heatOff", "service": { "serviceDefinition": "heating-switch" },
        "next": ["confirm"], "input": { "on": false } },
      { "name": "confirm", "service": { "serviceDefinition": "temperature-read" },
        "dep": ["heatOn", "heatOff"] }
    ] } ]
}
\end{lstlisting}


\begin{lstlisting}[caption={Example response body}]
{ "recipe": { "id": 7, "name": "temperature-branching", "version": "1.0.0",
    "status": "ACTIVE", "registeredBy": "WorkflowConfigurator",
    "recipe": { ... }, "createdAt": "2026-11-21T09:12:44Z" },
  "validation": { "recipe": "temperature-branching@1.0.0", "verdict": "ACCEPT", "findings": [ ... ] },
  "created": true }
\end{lstlisting}

\osubsection{query}{GET}{/workflowmanager/recipe/recipes}{RecipeFilter}{RecipeList}
Lists registered recipes in registration order, optionally only the versions of one recipe.

{\small
\begin{longtable}{|p{3.0cm}|p{1.4cm}|p{3.2cm}|p{8.2cm}|}

\hline
\rowcolor{gray!33} Parameter & In & Type & Description \\ \hline
\endhead
\fparam{name} & query & \pref{Name} & Optional recipe name. \\ \hline
\fparam{summary} & query & \pref{Boolean} & Optional, default false. When true, each entry leaves out \texttt{recipe}, so the list stays small however large the recipes are; \oref{get} returns one document. \\ \hline
\end{longtable}
}


{\small
\begin{longtable}{|p{2.0cm}|p{13.8cm}|}

\hline
\rowcolor{gray!33} Status & When \\ \hline
\endhead
Success & 200 OK. \\ \hline
401 & No valid identity token, or the token is bound to another certificate (\texttt{AUTH}). \\ \hline
403 & The consumer has no ConsumerAuthorization grant for this operation (\texttt{FORBIDDEN}). \\ \hline
\end{longtable}
}

\osubsection{get}{GET}{/workflowmanager/recipe/recipes/\{name\}/\{version\}}{RecipeKey}{RecipeRecord}
Returns one registered recipe version, active or retired.

{\small
\begin{longtable}{|p{3.0cm}|p{1.4cm}|p{3.2cm}|p{8.2cm}|}

\hline
\rowcolor{gray!33} Parameter & In & Type & Description \\ \hline
\endhead
\fparam{name} & path & \pref{Name} & Recipe name. \\ \hline
\fparam{version} & path & \pref{Version} & Recipe version. \\ \hline
\end{longtable}
}


{\small
\begin{longtable}{|p{2.0cm}|p{13.8cm}|}

\hline
\rowcolor{gray!33} Status & When \\ \hline
\endhead
Success & 200 OK. \\ \hline
404 & The referenced object does not exist (\texttt{DATA\_NOT\_FOUND}). \\ \hline
401 & No valid identity token, or the token is bound to another certificate (\texttt{AUTH}). \\ \hline
403 & The consumer has no ConsumerAuthorization grant for this operation (\texttt{FORBIDDEN}). \\ \hline
\end{longtable}
}

\osubsection{retire}{DELETE}{/workflowmanager/recipe/recipes/\{name\}/\{version\}}{RecipeKey}{RecipeRecord}
Marks a recipe version \texttt{RETIRED}. Running sessions continue; new sessions are refused.

{\small
\begin{longtable}{|p{3.0cm}|p{1.4cm}|p{3.2cm}|p{8.2cm}|}

\hline
\rowcolor{gray!33} Parameter & In & Type & Description \\ \hline
\endhead
\fparam{name} & path & \pref{Name} & Recipe name. \\ \hline
\fparam{version} & path & \pref{Version} & Recipe version. \\ \hline
\end{longtable}
}


{\small
\begin{longtable}{|p{2.0cm}|p{13.8cm}|}

\hline
\rowcolor{gray!33} Status & When \\ \hline
\endhead
Success & 200 OK with the retired record. \\ \hline
404 & The referenced object does not exist (\texttt{DATA\_NOT\_FOUND}). \\ \hline
401 & No valid identity token, or the token is bound to another certificate (\texttt{AUTH}). \\ \hline
403 & The consumer has no ConsumerAuthorization grant for this operation (\texttt{FORBIDDEN}). \\ \hline
\end{longtable}
}

\osubsection{compile}{POST}{/workflowmanager/recipe/compile}{RecipeKey}{ChoreographerPlan}
Compiles an active recipe version (the latest when \texttt{version} is absent) to a Choreographer plan request.

{\small
\begin{longtable}{|p{2.0cm}|p{13.8cm}|}

\hline
\rowcolor{gray!33} Status & When \\ \hline
\endhead
Success & 200 OK. \\ \hline
404 & The referenced object does not exist (\texttt{DATA\_NOT\_FOUND}). \\ \hline
409 & The recipe version is retired. \\ \hline
401 & No valid identity token, or the token is bound to another certificate (\texttt{AUTH}). \\ \hline
403 & The consumer has no ConsumerAuthorization grant for this operation (\texttt{FORBIDDEN}). \\ \hline
\end{longtable}
}


\begin{lstlisting}[caption={Example request body}]
{ "name": "temperature-control", "version": "1.0.0" }
\end{lstlisting}


\begin{lstlisting}[caption={Example response body}]
{ "name": "temperature-control", "firstActionName": "control",
  "actions": [ { "name": "control", "firstStepNames": ["read"],
    "steps": [ { "name": "read", "quantity": 1,
      "serviceRequirement": { "serviceDefinition": "temperature-read" },
      "localCloudOnly": true, "nextStepNames": ["evaluate"] }, ... ] } ],
  "lossySteps": ["evaluate", "switch"] }
\end{lstlisting}


\newpage
\section{Data Models}
\label{sec:model}
The JSON objects exchanged by the operations, in alphabetic order. Field types refer to other data models or to the
primitives of Section~\ref{sec:model:primitives}.

\msubsection{struct}{Action}
An action $\alpha$: a group of steps run in one work centre.

{\small
\begin{longtable}{|p{3.7cm}|p{4.2cm}|p{1.6cm}|p{6.3cm}|}

\hline
\rowcolor{gray!33} Field & Type & Mandatory & Description \\ \hline
\endhead
\texttt{name} & \pref{Name} & yes & Action name. \\ \hline
\texttt{next} & \pref{Name} & no & The action that becomes active when this one finishes; absent for the last action. \\ \hline
\texttt{workCentre} & \pref{String} & no & Informative work-centre identifier. \\ \hline
\texttt{dispatch} & \mref{DispatchMode} & no & Dispatch mode of this action's steps. \\ \hline
\texttt{steps} & \pref{List}$<$\mref{Step}$>$ & yes & The steps $S_\alpha$. \\ \hline
\end{longtable}
}

\msubsection{struct}{ChoreographerPlan}
A recipe compiled to the Choreographer plan request (\texttt{PlanRequestDTO}), for option A of the coexistence with the Choreographer.

{\small
\begin{longtable}{|p{3.7cm}|p{4.2cm}|p{1.6cm}|p{6.3cm}|}

\hline
\rowcolor{gray!33} Field & Type & Mandatory & Description \\ \hline
\endhead
\texttt{name} & \pref{Name} & yes & Plan name: the recipe name. \\ \hline
\texttt{firstActionName} & \pref{Name} & yes & First action. \\ \hline
\texttt{actions} & \pref{List}$<$\mref{PlanAction}$>$ & yes & Actions. \\ \hline
\texttt{lossySteps} & \pref{List}$<$\pref{String}$>$ & no & Informational: steps whose input, guards, dependencies or policy the plan cannot express. Not part of the Choreographer request. \\ \hline
\end{longtable}
}

\msubsection{enum}{DispatchMode}
How a step reaches an executor.
 Encoded as a JSON string.
{\small
\begin{longtable}{|p{3.9cm}|p{11.9cm}|}

\hline
\rowcolor{gray!33} Value & Meaning \\ \hline
\endhead
\texttt{push} & The WorkflowManager selects an executor and calls \texttt{workflowExecution.start} (default). \\ \hline
\texttt{pull} & The WorkflowManager publishes the work item; an executor claims it with \texttt{workflowWorkItem.claim}. \\ \hline
\texttt{event} & Reserved for EventHandler notification; currently handled as \texttt{pull}. \\ \hline
\end{longtable}
}

\msubsection{struct}{ErrorMessage}
Error response of every operation (common library \texttt{ErrorMessageDTO}).

{\small
\begin{longtable}{|p{3.7cm}|p{4.2cm}|p{1.6cm}|p{6.3cm}|}

\hline
\rowcolor{gray!33} Field & Type & Mandatory & Description \\ \hline
\endhead
\texttt{errorMessage} & \pref{String} & yes & Explanation. \\ \hline
\texttt{errorCode} & \pref{Integer} & yes & HTTP status code. \\ \hline
\texttt{exceptionType} & \pref{String} & yes & \texttt{INVALID\_PARAMETER}, \texttt{AUTH}, \texttt{FORBIDDEN}, \texttt{DATA\_NOT\_FOUND}, \texttt{EXTERNAL\_SERVER\_ERROR}, \ldots \\ \hline
\texttt{origin} & \pref{String} & no & The operation that failed, e.g. \texttt{POST /workflowmanager/work-item/report}. \\ \hline
\end{longtable}
}

\msubsection{struct}{ExecutorRequirement}
$\eta(s)$: constraints on the executor that runs a step.

{\small
\begin{longtable}{|p{3.7cm}|p{4.2cm}|p{1.6cm}|p{6.3cm}|}

\hline
\rowcolor{gray!33} Field & Type & Mandatory & Description \\ \hline
\endhead
\texttt{capability} & \pref{String} & no & A capability the executor must declare in \texttt{workflowCapabilities}. \\ \hline
\texttt{workstation} & \pref{String} & no & The executor's workstation. \\ \hline
\texttt{system} & \pref{SystemName} & no & A named executor system. \\ \hline
\end{longtable}
}

\msubsection{struct}{Finding}
One validator finding.

{\small
\begin{longtable}{|p{3.7cm}|p{4.2cm}|p{1.6cm}|p{6.3cm}|}

\hline
\rowcolor{gray!33} Field & Type & Mandatory & Description \\ \hline
\endhead
\texttt{code} & \pref{String} & yes & Check code, e.g. \texttt{D03}; the codes are listed in the \texttt{workflowRecipe} SD. \\ \hline
\texttt{outcome} & \mref{Verdict} & yes & Outcome of this finding. \\ \hline
\texttt{basis} & \pref{String} & yes & The rule the check enforces, e.g. \texttt{W1}, \texttt{V9}, \texttt{schema}. \\ \hline
\texttt{location} & \pref{String} & yes & Where the finding applies, e.g. \texttt{step evaluate}. \\ \hline
\texttt{message} & \pref{String} & yes & Human-readable explanation. \\ \hline
\end{longtable}
}

\msubsection{enum}{LeaseMode}
Role of a step for a shared resource.
 Encoded as a JSON string.
{\small
\begin{longtable}{|p{3.9cm}|p{11.9cm}|}

\hline
\rowcolor{gray!33} Value & Meaning \\ \hline
\endhead
\texttt{reserve} & Reserves the resource; its result carries the lease token. \\ \hline
\texttt{use} & Uses the resource; the executor renews the lease immediately before the invocation. \\ \hline
\texttt{release} & Releases the resource. \\ \hline
\end{longtable}
}

\msubsection{struct}{LeasePolicy}
The role of a step for a shared resource.

{\small
\begin{longtable}{|p{3.7cm}|p{4.2cm}|p{1.6cm}|p{6.3cm}|}

\hline
\rowcolor{gray!33} Field & Type & Mandatory & Description \\ \hline
\endhead
\texttt{resource} & \pref{Name} & yes & A declared resource. \\ \hline
\texttt{mode} & \mref{LeaseMode} & yes & \texttt{reserve}, \texttt{use} or \texttt{release}. \\ \hline
\end{longtable}
}

\msubsection{enum}{OnError}
Reaction to an \texttt{ERROR} report.
 Encoded as a JSON string.
{\small
\begin{longtable}{|p{3.9cm}|p{11.9cm}|}

\hline
\rowcolor{gray!33} Value & Meaning \\ \hline
\endhead
\texttt{rebind} & Exclude the executor for this session and bind the step again. \\ \hline
\texttt{retry} & Keep the executor: after \texttt{backoffSeconds} dispatch the step to the same executor again (extension x-retry; for transient provider-side failures, the provider must accept a repeated call). \\ \hline
\texttt{abort} & Fail the session. \\ \hline
\end{longtable}
}

\msubsection{struct}{Parameter}
A declared session parameter.

{\small
\begin{longtable}{|p{3.7cm}|p{4.2cm}|p{1.6cm}|p{6.3cm}|}

\hline
\rowcolor{gray!33} Field & Type & Mandatory & Description \\ \hline
\endhead
\texttt{name} & \pref{Name} & yes & Parameter name, referenced as \texttt{\$\{param.name\}}. \\ \hline
\texttt{type} & \pref{String} & no & Informative type: \texttt{string}, \texttt{number}, \texttt{boolean} or \texttt{object}. \\ \hline
\texttt{required} & \pref{Boolean} & no & Whether \texttt{workflowInstance.start} refuses a session without it. \\ \hline
\end{longtable}
}

\msubsection{struct}{PlanAction}
\texttt{ActionRequestDTO} of the Choreographer.

{\small
\begin{longtable}{|p{3.7cm}|p{4.2cm}|p{1.6cm}|p{6.3cm}|}

\hline
\rowcolor{gray!33} Field & Type & Mandatory & Description \\ \hline
\endhead
\texttt{name} & \pref{Name} & yes & Action name. \\ \hline
\texttt{nextActionName} & \pref{Name} & no & Next action. \\ \hline
\texttt{firstStepNames} & \pref{List}$<$\pref{Name}$>$ & yes & Steps without predecessors. \\ \hline
\texttt{steps} & \pref{List}$<$\mref{PlanStep}$>$ & yes & Steps. \\ \hline
\end{longtable}
}

\msubsection{struct}{PlanServiceRequirement}
Subset of the orchestration service requirement.

{\small
\begin{longtable}{|p{3.7cm}|p{4.2cm}|p{1.6cm}|p{6.3cm}|}

\hline
\rowcolor{gray!33} Field & Type & Mandatory & Description \\ \hline
\endhead
\texttt{serviceDefinition} & \pref{ServiceDefinition} & yes & Service. \\ \hline
\texttt{versions} & \pref{List}$<$\pref{String}$>$ & no & Accepted versions. \\ \hline
\texttt{interfaceTemplateNames} & \pref{List}$<$\pref{String}$>$ & no & Accepted interface templates. \\ \hline
\end{longtable}
}

\msubsection{struct}{PlanStep}
\texttt{StepRequestDTO} of the Choreographer.

{\small
\begin{longtable}{|p{3.7cm}|p{4.2cm}|p{1.6cm}|p{6.3cm}|}

\hline
\rowcolor{gray!33} Field & Type & Mandatory & Description \\ \hline
\endhead
\texttt{name} & \pref{Name} & yes & Step name. \\ \hline
\texttt{quantity} & \pref{Integer} & no & Always 1. \\ \hline
\texttt{serviceRequirement} & \mref{PlanServiceRequirement} & yes & Service of the step. \\ \hline
\texttt{localCloudOnly} & \pref{Boolean} & no & From the step policy. \\ \hline
\texttt{staticParameters} & \pref{Map}$<$\pref{String}, \pref{String}$>$ & no & The step's \texttt{params} as strings. \\ \hline
\texttt{nextStepNames} & \pref{List}$<$\pref{Name}$>$ & no & Successors. \\ \hline
\end{longtable}
}

\msubsection{struct}{Policy}
Execution policy of a step. Every field other than \texttt{lease} and \texttt{localCloudOnly} is an extension outside the proven model.

{\small
\begin{longtable}{|p{3.7cm}|p{4.2cm}|p{1.6cm}|p{6.3cm}|}

\hline
\rowcolor{gray!33} Field & Type & Mandatory & Description \\ \hline
\endhead
\texttt{quantity} & \pref{Integer} & no & Number of executions; only 1 is accepted (check M01). \\ \hline
\texttt{timeoutSeconds} & \pref{Duration} & no & Extension x-timeout: the step fails when it runs longer. \\ \hline
\texttt{nominalDurationSeconds} & \pref{Duration} & no & Expected duration, used to check lease durations (check L04). \\ \hline
\texttt{retry} & \mref{Retry} & no & Extension x-rebind: behaviour on \texttt{ERROR}. \\ \hline
\texttt{lease} & \mref{LeasePolicy} & no & The step's role for a shared resource. \\ \hline
\texttt{localCloudOnly} & \pref{Boolean} & no & Whether orchestration may return providers of other clouds (default true). \\ \hline
\end{longtable}
}

\msubsection{struct}{Recipe}
A workflow recipe $W$~\cite{varga2026}: named, versioned and immutable once registered. Actions run in sequence (\texttt{firstAction}, then each action's \texttt{next}); the steps of one action form a directed acyclic graph.

{\small
\begin{longtable}{|p{3.7cm}|p{4.2cm}|p{1.6cm}|p{6.3cm}|}

\hline
\rowcolor{gray!33} Field & Type & Mandatory & Description \\ \hline
\endhead
\texttt{name} & \pref{Name} & yes & Recipe name. \\ \hline
\texttt{version} & \pref{Version} & yes & Recipe version; a registered name and version never changes content. \\ \hline
\texttt{description} & \pref{String} & no & Free text. \\ \hline
\texttt{firstAction} & \pref{Name} & yes & The action $\alpha_0$ that is active when a session starts. \\ \hline
\texttt{parameters} & \pref{List}$<$\mref{Parameter}$>$ & no & Declared session parameters $\rho$. \\ \hline
\texttt{resources} & \pref{List}$<$\mref{Resource}$>$ & no & Shared resources (tools) used under leases. \\ \hline
\texttt{dispatch} & \mref{DispatchMode} & no & Default dispatch mode of the recipe; overridden per action or per session. \\ \hline
\texttt{actions} & \pref{List}$<$\mref{Action}$>$ & yes & The actions $A$ with their steps. \\ \hline
\end{longtable}
}

\msubsection{struct}{RecipeFilter}
Optional filter of the recipe query.

{\small
\begin{longtable}{|p{3.7cm}|p{4.2cm}|p{1.6cm}|p{6.3cm}|}

\hline
\rowcolor{gray!33} Field & Type & Mandatory & Description \\ \hline
\endhead
\texttt{name} & \pref{Name} & no & Only versions of this recipe. \\ \hline
\texttt{summary} & \pref{Boolean} & no & Leave out the recipe documents (default false). \\ \hline
\end{longtable}
}

\msubsection{struct}{RecipeKey}
Identifies one registered recipe version.

{\small
\begin{longtable}{|p{3.7cm}|p{4.2cm}|p{1.6cm}|p{6.3cm}|}

\hline
\rowcolor{gray!33} Field & Type & Mandatory & Description \\ \hline
\endhead
\texttt{name} & \pref{Name} & yes & Recipe name. \\ \hline
\texttt{version} & \pref{Version} & yes & Recipe version; for \texttt{compile} the latest active version when absent. \\ \hline
\end{longtable}
}

\msubsection{struct}{RecipeList}
Registered recipes in registration order.

{\small
\begin{longtable}{|p{3.7cm}|p{4.2cm}|p{1.6cm}|p{6.3cm}|}

\hline
\rowcolor{gray!33} Field & Type & Mandatory & Description \\ \hline
\endhead
\texttt{entries} & \pref{List}$<$\mref{RecipeRecord}$>$ & yes & The recipes. \\ \hline
\texttt{count} & \pref{Integer} & yes & Number of entries. \\ \hline
\end{longtable}
}

\msubsection{struct}{RecipeRecord}
A registered recipe as stored by the WorkflowManager.

{\small
\begin{longtable}{|p{3.7cm}|p{4.2cm}|p{1.6cm}|p{6.3cm}|}

\hline
\rowcolor{gray!33} Field & Type & Mandatory & Description \\ \hline
\endhead
\texttt{id} & \pref{Integer} & yes & Database identifier. \\ \hline
\texttt{name} & \pref{Name} & yes & Recipe name. \\ \hline
\texttt{version} & \pref{Version} & yes & Recipe version. \\ \hline
\texttt{status} & \mref{RecipeStatus} & yes & \texttt{ACTIVE} or \texttt{RETIRED}. \\ \hline
\texttt{registeredBy} & \pref{SystemName} & yes & The system that registered the recipe. \\ \hline
\texttt{recipe} & \mref{Recipe} & yes & The canonical recipe document; null in a \texttt{workflowRecipe.query} with \texttt{summary=true}. \\ \hline
\texttt{createdAt} & \pref{DateTime} & yes & Registration time. \\ \hline
\end{longtable}
}

\msubsection{struct}{RecipeRegistration}
Result of a registration.

{\small
\begin{longtable}{|p{3.7cm}|p{4.2cm}|p{1.6cm}|p{6.3cm}|}

\hline
\rowcolor{gray!33} Field & Type & Mandatory & Description \\ \hline
\endhead
\texttt{recipe} & \mref{RecipeRecord} & no & The stored recipe; null when the recipe was not accepted. \\ \hline
\texttt{validation} & \mref{ValidationReport} & yes & The validation report. \\ \hline
\texttt{created} & \pref{Boolean} & yes & True when a new version was stored; false when identical content already existed. \\ \hline
\end{longtable}
}

\msubsection{enum}{RecipeStatus}
Life-cycle state of a registered recipe.
 Encoded as a JSON string.
{\small
\begin{longtable}{|p{3.9cm}|p{11.9cm}|}

\hline
\rowcolor{gray!33} Value & Meaning \\ \hline
\endhead
\texttt{ACTIVE} & Sessions may be started. \\ \hline
\texttt{RETIRED} & Kept for traceability; new sessions are refused with 409. \\ \hline
\end{longtable}
}

\msubsection{struct}{Resource}
A shared resource that steps reserve, use and release under a lease (Proposition~1 of~\cite{varga2026}, lease variant).

{\small
\begin{longtable}{|p{3.7cm}|p{4.2cm}|p{1.6cm}|p{6.3cm}|}

\hline
\rowcolor{gray!33} Field & Type & Mandatory & Description \\ \hline
\endhead
\texttt{name} & \pref{Name} & yes & Resource name. \\ \hline
\texttt{leaseDurationSeconds} & \pref{Duration} & no & Lease duration granted by \texttt{workflowLease}; must exceed the longest tool-using step. \\ \hline
\texttt{provider} & \pref{SystemName} & no & Lease provider; the WorkflowManager when absent. \\ \hline
\end{longtable}
}

\msubsection{struct}{Retry}
Extensions x-rebind and x-retry.

{\small
\begin{longtable}{|p{3.7cm}|p{4.2cm}|p{1.6cm}|p{6.3cm}|}

\hline
\rowcolor{gray!33} Field & Type & Mandatory & Description \\ \hline
\endhead
\texttt{maxAttempts} & \pref{Integer} & no & Maximum attempts of the step in one session. \\ \hline
\texttt{onError} & \mref{OnError} & no & What an \texttt{ERROR} report causes. \\ \hline
\texttt{backoffSeconds} & \pref{Duration} & no & Pause before a same-executor retry (\texttt{onError} = \texttt{retry}); 0--600, default 1. \\ \hline
\end{longtable}
}

\msubsection{struct}{ServiceRequirement}
Requirements on the provider of the step's service, passed to orchestration.

{\small
\begin{longtable}{|p{3.7cm}|p{4.2cm}|p{1.6cm}|p{6.3cm}|}

\hline
\rowcolor{gray!33} Field & Type & Mandatory & Description \\ \hline
\endhead
\texttt{serviceDefinition} & \pref{ServiceDefinition} & yes & The service to consume. \\ \hline
\texttt{interfaces} & \pref{List}$<$\pref{String}$>$ & no & Accepted interface template names, e.g. \texttt{generic\_https}. \\ \hline
\texttt{minVersion} & \pref{String} & no & Lowest accepted service version. \\ \hline
\texttt{maxVersion} & \pref{String} & no & Highest accepted service version. \\ \hline
\texttt{metadata} & \pref{Map}$<$\pref{String}, \pref{JsonValue}$>$ & no & Metadata requirements on the service instance. \\ \hline
\texttt{qos} & \pref{Map}$<$\pref{String}, \pref{JsonValue}$>$ & no & QoS requirements forwarded to orchestration. \\ \hline
\texttt{requiresResults} & \pref{List}$<$\pref{Name}$>$ & no & Steps whose results the provider needs; each must be in $dep(s)$ (check D04). \\ \hline
\end{longtable}
}

\msubsection{struct}{Step}
A step $s$: one service exchange.

{\small
\begin{longtable}{|p{3.7cm}|p{4.2cm}|p{1.6cm}|p{6.3cm}|}

\hline
\rowcolor{gray!33} Field & Type & Mandatory & Description \\ \hline
\endhead
\texttt{name} & \pref{Name} & yes & Step name, unique in the recipe. \\ \hline
\texttt{service} & \mref{ServiceRequirement} & yes & $\theta(s)$: the service the step consumes. \\ \hline
\texttt{next} & \pref{List}$<$\pref{Name}$>$ & no & Successors in the same action (the transitions $s \rightarrow s'$). \\ \hline
\texttt{dep} & \pref{List}$<$\pref{Name}$>$ & no & $dep(s)$: steps whose results the step reads. Each must run before $s$. \\ \hline
\texttt{executor} & \mref{ExecutorRequirement} & no & $\eta(s)$: which executor may run the step; absent means any executor declaring the service. \\ \hline
\texttt{params} & \pref{Map}$<$\pref{String}, \pref{JsonValue}$>$ & no & Static parameters passed unchanged to the provider. \\ \hline
\texttt{input} & \pref{Template} & no & $\kappa_s$: the request body sent to the provider, composed from $\rho$ and $\xi|dep(s)$. \\ \hline
\texttt{policy} & \mref{Policy} & no & Execution policy (quantity, timeout, retry, lease, local cloud). \\ \hline
\texttt{guards} & \pref{Map}$<$\pref{Name}, \pref{Expression}$>$ & no & Conditions on the transitions to successors (extension x-guard). A successor without a guard is always taken. \\ \hline
\end{longtable}
}

\msubsection{struct}{ValidationReport}
The result of validating a recipe. The verdict is the worst outcome of the findings (REJECT $>$ CLARIFY $>$ REVISE $>$ ACCEPT).

{\small
\begin{longtable}{|p{3.7cm}|p{4.2cm}|p{1.6cm}|p{6.3cm}|}

\hline
\rowcolor{gray!33} Field & Type & Mandatory & Description \\ \hline
\endhead
\texttt{recipe} & \pref{String} & yes & \texttt{name@version} of the validated recipe. \\ \hline
\texttt{verdict} & \mref{Verdict} & yes & Overall verdict; only \texttt{ACCEPT} can be registered. \\ \hline
\texttt{findings} & \pref{List}$<$\mref{Finding}$>$ & yes & Individual findings; empty when the recipe has none. \\ \hline
\end{longtable}
}

\msubsection{enum}{Verdict}
Validator outcome.
 Encoded as a JSON string.
{\small
\begin{longtable}{|p{3.9cm}|p{11.9cm}|}

\hline
\rowcolor{gray!33} Value & Meaning \\ \hline
\endhead
\texttt{ACCEPT} & The recipe can be registered (also used for informational findings). \\ \hline
\texttt{REVISE} & The recipe must be changed. \\ \hline
\texttt{CLARIFY} & The recipe refers to something that cannot be confirmed (undeclared parameter, unknown service or executor). \\ \hline
\texttt{REJECT} & The recipe cannot be used at all. \\ \hline
\end{longtable}
}


\subsection{Primitives}
\label{sec:model:primitives}
The abstract primitive types of the SD are represented in JSON~\cite{rfc8259} as follows.
\subsubsection*{alias \pdef{Boolean}}
JSON \texttt{true} or \texttt{false}.

\subsubsection*{alias \pdef{DateTime}}
JSON string in the RFC 3339 format~\cite{rfc3339}, e.g. \texttt{2026-11-23T10:15:30Z}.

\subsubsection*{alias \pdef{Duration}}
JSON number of seconds.

\subsubsection*{alias \pdef{Expression}}
JSON string in the guard grammar.

\subsubsection*{alias \pdef{Integer}}
JSON number without fraction.

\subsubsection*{alias \pdef{JsonValue}}
Any JSON value~\cite{rfc8259}.

\subsubsection*{alias \pdef{List}}
JSON array.

\subsubsection*{alias \pdef{Map}}
JSON object whose member names are the keys.

\subsubsection*{alias \pdef{Name}}
JSON string matching \texttt{\^{}[A-Za-z][A-Za-z0-9\_-]\{0,62\}\$}.

\subsubsection*{alias \pdef{ServiceDefinition}}
JSON string.

\subsubsection*{alias \pdef{String}}
JSON string (UTF-8).

\subsubsection*{alias \pdef{SystemName}}
JSON string.

\subsubsection*{alias \pdef{Template}}
Any JSON value; references are JSON strings.

\subsubsection*{alias \pdef{Version}}
JSON string matching \texttt{\^{}\textbackslash{}d+\textbackslash{}.\textbackslash{}d+\textbackslash{}.\textbackslash{}d+\$}.



\section{Security}
\label{sec:security}
The interface is offered only over TLS~\cite{rfc8446} (TLS 1.3 or 1.2) with mutual authentication: the consumer must
present an Arrowhead certificate issued under the local cloud's certificate authority, and certificates on the
cloud's CRL are refused. The service is registered with interface template \texttt{generic\_https} and interface
policy \texttt{CERT\_AUTH}.

Every request must carry the consumer's identity token, obtained from the Authentication system and bound to the
consumer's certificate:
\begin{lstlisting}
Authorization: Bearer IDENTITY-TOKEN//<identity token>
\end{lstlisting}
A token issued to another system's certificate is refused (401). The provider then checks a ConsumerAuthorization
grant for the operation (\texttt{authorization.verify}, \texttt{service.authorization.enabled=true}); a missing grant is
refused with 403. Each operation is authorized separately (\texttt{workflowRecipe.validate}, \ldots). The WorkflowConfigurator needs \texttt{validate}, \texttt{register}, \texttt{query} and \texttt{get}.


\section{Errors}
\label{sec:errors}
Every error response carries an \mref{ErrorMessage} body with content type \texttt{application/json}, for example:
\begin{lstlisting}
{ "errorMessage": "Work item 5d2c0b1e-... is held by another executor",
  "errorCode": 403, "exceptionType": "FORBIDDEN",
  "origin": "POST /workflowmanager/work-item/report" }
\end{lstlisting}
Status codes follow HTTP semantics~\cite{rfc9110}: 400 invalid parameter, 401 authentication, 403 forbidden, 404 not
found, 409 conflict with the current state, 500 internal error and 503 failure of an external system.


\section{Implementations}
\label{sec:implementations}
\begin{itemize}
\item \textbf{Java (reference).} \texttt{Eclipse-Arrowhead-java/ah5-Workflow-java}, module \texttt{workflow-manager}. Paths, methods and service
operations are defined once in \texttt{workflow-model} (\texttt{wire/WorkflowServices.java}, \texttt{wire/ServiceOperations.java}) and checked by the provider-side authorization filter.
\item \textbf{Go.} Not implemented: \texttt{Eclipse-Arrowhead-go} has no workflow systems. A Go implementation must
provide exactly the paths, methods, JSON field names and status codes of this IDD to be interoperable.
\item \textbf{Rust.} Not implemented: \texttt{Eclipse-Arrowhead-rust} has no workflow systems. The same requirement applies.
\end{itemize}
Until the Go and Rust implementations exist, the workflow systems are \emph{not} interoperable across the three
implementations; only the Java systems interoperate with Java, Go and Rust core systems through the common Arrowhead
interfaces (ServiceRegistry, Authentication, ConsumerAuthorization, DynamicServiceOrchestration).


\newpage
\bibliographystyle{IEEEtran}
\bibliography{bibliography}

\newpage

\section{Revision History}
\subsection{Amendments}

\noindent\begin{tabularx}{\textwidth}{| p{1cm} | p{2.2cm} | p{1.6cm} | X | p{3.6cm} |} \hline
\rowcolor{gray!33} No. & Date & Version & Subject of Amendments & Author \\ \hline
1 & 2026-11-23 & \arrowversion & Initial edition: interface design description of workflowRecipe over HTTP/TLS/JSON, written from the Java reference implementation (ah5-Workflow-java). & Documentation maintenance \\ \hline
\end{tabularx}

\subsection{Quality Assurance}

\noindent\begin{tabularx}{\textwidth}{| p{1cm} | p{2.2cm} | p{1.6cm} | X |} \hline
\rowcolor{gray!33} No. & Date & Version & Approved by \\ \hline
1 & & \arrowversion & Pending review \\ \hline
\end{tabularx}

\end{document}
